\documentclass[11pt]{article}

\usepackage[T1]{fontenc}
\usepackage{lmodern}
\usepackage[margin=1in]{geometry}
\usepackage{amsmath,amssymb,mathtools}
\usepackage{booktabs}
\usepackage{microtype}
\usepackage[numbers,sort&compress]{natbib}
\usepackage[hidelinks]{hyperref}
\hypersetup{
  pdftitle={A two-pineapple envelope for split-core threshold random walks},
  pdfauthor={Xiaojian Zeng}
}

\newcommand{\K}{\mathcal{K}}
\newcommand{\N}{\mathbb{N}}
\newcommand{\one}{\mathbf{1}}

\title{A two-pineapple envelope for split-core threshold random walks}
\author{Xiaojian Zeng\\
\small PBC School of Finance, Tsinghua University, Beijing 100190, China\\
\small \texttt{xiaojianzeng@foxmail.com}}
\date{}

\begin{document}
\maketitle

\begin{abstract}
Kemeny's constant is a stationary-target mean first-passage time for a finite
random walk.  A conjecture of Breen et al. proposes that, among threshold graphs
of a fixed order, its maximum is attained by a pineapple graph.  We prove an
exact partial result for two infinite split-core families.  For positive integers
$s,q,t$ and $n=s+q+t+3$, each threshold code
$01\,0^t1^s0^q1$ and $0\,1^s0^t1\,0^q1$ has Kemeny's constant strictly below
the larger of the two pineapple values with core parameters $s+1$ and
$s+t+1$.  The proof reduces the general construction-code formula to four
rational polynomials and uses exact positive-coefficient certificates.  We also
show that the tempting operation of filling the interior zero block is not
monotone: it can increase, decrease, or leave Kemeny's constant unchanged.
Thus the two-pineapple comparison is essential.  The global pineapple conjecture
and other split-core placements remain unresolved.
\end{abstract}

\noindent\textbf{Keywords:} Kemeny's constant; threshold graph; random walk;
pineapple graph; exact symbolic computation; extremal graph theory

\section{Introduction}

Let a simple random walk move one edge per unit time on a finite connected
graph.  Kemeny's constant $\K$ is the expected time to reach a target drawn
from the stationary distribution, and it does not depend on the initial
vertex.  It is simultaneously a transport statistic and a spectral quantity.
The way in which graph structure controls $\K$ is subtle: adding edges changes
degrees, transition probabilities, and stationary weights at the same time.

Threshold graphs provide a sharply structured setting in which this extremal
question remains nontrivial.  Breen et al.~\cite{BreenEtAl2025} derived an exact
construction-code formula, studied several related walk parameters, and
conjectured that a pineapple graph maximizes $\K$ among all threshold graphs of
fixed order.  Their openly available preprint reports computational verification
through order $30$, but the conjecture is universal.

This paper isolates two natural ways to split a dense core by a nonempty block
of isolated-at-insertion vertices.  Neither family can beat the pineapple
envelope.  More precisely, each member is strictly dominated by one of two
pineapples of the same order.  Which pineapple wins depends on the parameters;
this dependence cannot be removed by a simple edge-filling monotonicity argument.
The result excludes two three-parameter families from a search for a global
counterexample, but it does not settle the full conjecture.

All calculations below are exact.  The source formula was transcribed from
arXiv:2310.08552v1, the version that was available for equation-level checking.
The corresponding 2025 version of record was not available for a line-by-line
comparison, so statements about proximity to prior work are intentionally
bounded to the sources described in Section~\ref{sec:related}.

\section{Threshold-walk model and the pineapple bound}
\label{sec:model}

A connected threshold graph on $n\geq 3$ vertices is encoded here by
$c=(c_1,\ldots,c_n)\in\{0,1\}^n$ with $c_1=0$ and $c_n=1$.  A vertex coded by
$1$ is joined to every earlier vertex, while a vertex coded by $0$ is joined
to none of them.  Equivalently,
\begin{equation}
 A_{ii}=0,
 \qquad A_{ij}=c_{\max(i,j)}\quad(i\neq j).
 \label{eq:adjacency}
\end{equation}
Let $d_i=\sum_j A_{ij}$, $D=\operatorname{diag}(d_i)$, and
$m=\frac12\sum_i d_i$.  The non-lazy discrete-time walk has transition matrix
$T=D^{-1}A$ and stationary distribution $\pi_i=d_i/(2m)$.  If
\begin{equation}
 H_{ij}=\mathbb E_i\inf\{k\geq0:X_k=j\},\qquad H_{ii}=0,
\end{equation}
then
\begin{equation}
 \K(G)=\sum_j\pi_jH_{ij}
 =\operatorname{Tr}\!\left(I-T+\one\pi^{\mathsf T}\right)^{-1}-1
 =\sum_{k=2}^n\frac{1}{1-\lambda_k(T)}.
 \label{eq:kemeny}
\end{equation}
The auxiliary first-hitting problem does not alter the original chain.  In
particular, periodic graphs are permitted and every step retains unit duration.

For $0\leq r\leq n-2$, the pineapple graph $G_{n,r}$ has code
$0\,1^r0^{n-r-2}1$.  Its exact value is~\cite[Lemma 3.14]{BreenEtAl2025}
\begin{equation}
 F_n(r)=n-4+\frac{2}{r+2}
 +\frac{(n-1)(2r+3)}{2n+r^2+r-2}.
 \label{eq:pineapple}
\end{equation}
Write $B_n=\max_{0\leq r\leq n-2}F_n(r)$, with the maximum taken over integer
$r$.  The global conjecture is $\K(G(c))\leq B_n$ for every legal code $c$.

For completeness, the source reduction used in our calculation can be stated
compactly.  For $1\leq i\leq n-1$, define vectors $w_i,z_i\in\mathbb R^n$ by
\begin{align}
 w_i&=(0,2,4,\ldots,2(i-1),-i(i-1),0,\ldots,0),\\
 z_i&=(0,\ldots,0,i+1,1,\ldots,1),
\end{align}
where the distinguished entry is in coordinate $i+1$.  Theorem 3.8 of
Ref.~\cite{BreenEtAl2025} gives
\begin{equation}
 \K(c)=n-1-\sum_{i=1}^{n-1}\frac{c_{i+1}}{z_i^{\mathsf T}c}
 +\sum_{i=1}^{n-1}
 \frac{(w_i^{\mathsf T}c)(2m-w_i^{\mathsf T}c)}
 {2m\,i(i+1)z_i^{\mathsf T}c}.
 \label{eq:sourceformula}
\end{equation}
This expression and Eq.~\eqref{eq:pineapple} were independently checked
against literal full-graph hitting-time calculations before the new comparison
was attempted.

\section{Exact two-pineapple domination}
\label{sec:theorem}

For $s,q,t\in\N$ define the legal six-run codes
\begin{equation}
 A_{s,q,t}=01\,0^t1^s0^q1,
 \qquad
 C_{s,q,t}=0\,1^s0^t1\,0^q1,
 \label{eq:families}
\end{equation}
both of order $n=s+q+t+3$.  We use $C$ for the second family to avoid
confusing the graph with the scalar pineapple envelope $B_n$.

\paragraph{Theorem 1 (split-core envelope).}
For every $s,q,t\in\N$,
\begin{equation}
 \K(A_{s,q,t}),\ \K(C_{s,q,t})
 <\max\{F_n(s+1),F_n(s+t+1)\}\leq B_n.
 \label{eq:main}
\end{equation}

\subsection{Block reduction}

On a constant-bit run $a\leq i\leq b$, both $z_i^{\mathsf T}c$ and
$w_i^{\mathsf T}c$ are constant.  If the run consists of ones, their values
are respectively
\begin{equation}
 b+1+\#\{\text{later ones}\},
 \qquad 2P_a-a(a-1),
\end{equation}
where $P_a$ is the sum of the indices of earlier ones.  On a zero run the
corresponding values are the number of later ones and $2P_a$.  The remaining
factor telescopes:
\begin{equation}
 \sum_{i=a}^{b}\frac{1}{i(i+1)}=\frac1a-\frac1{b+1}.
 \label{eq:telescope}
\end{equation}
Thus Eq.~\eqref{eq:sourceformula} reduces each code in Eq.~\eqref{eq:families}
to a rational function of $s,q,t$ without a matrix diagonalization.

Set
\begin{align}
 \Delta_X^h&=F_n(s+t+1)-\K(X),&
 \Delta_X^\ell&=F_n(s+1)-\K(X),
\end{align}
where $X=A_{s,q,t}$ or $C_{s,q,t}$.  Exact simplification gives
\begin{align}
 \Delta_A^h&=\frac{tH_A}{D_A^h}, &
 \Delta_A^\ell&=\frac{stL_A}{D_A^\ell},\label{eq:Agaps}\\
 \Delta_C^h&=\frac{tH_C}{D_C^h}, &
 \Delta_C^\ell&=\frac{tL_C}{D_C^\ell}.\label{eq:Cgaps}
\end{align}
Here $H_X,L_X$ are integer polynomials fixed by these normalized expressions.
Their complete expanded forms accompany this article as an exact
machine-readable certificate.  The denominators are
\begin{align*}
D_A^h={}&(s+1)(s+3)(s+t+3)
 (2q+s^2+2st+5s+2t+6)\\
&\times(2q+s^2+2st+5s+t^2+5t+6),\\
D_A^\ell={}&(s+1)(s+3)(s+t+3)
 (2q+s^2+5s+2t+6)\\
&\times(2q+s^2+2st+5s+2t+6),\\
D_C^h={}&(s+3)(s+t+3)(2q+s^2+5s+4t+6)\\
&\times(2q+s^2+2st+5s+t^2+5t+6),\\
D_C^\ell={}&(s+3)(s+t+3)(2q+s^2+5s+2t+6)\\
&\times(2q+s^2+5s+4t+6).
\end{align*}
Every displayed factor is positive on the stated domain.

\subsection{Positive-coefficient certificates}

For the first family, exact polynomial division yields
\begin{equation}
 L_A+\alpha_AH_A=R_A,
 \qquad
 \alpha_A=\frac{s^2+4s+1}{s^2t+3s^2+4st+12s+3t+7}>0.
 \label{eq:Acertificate}
\end{equation}
The denominator of $R_A$ is the denominator in $\alpha_A$, and its numerator
is an expanded polynomial with 49 nonzero monomials.  Every coefficient is a
strictly positive integer.  The full polynomial is retained verbatim in the
supplementary certificate rather than abbreviated by a numerical fit.
Consequently, if $H_A>0$, Eq.~\eqref{eq:Agaps} gives $\Delta_A^h>0$; if
$H_A\leq0$, Eq.~\eqref{eq:Acertificate} gives
$L_A=R_A-\alpha_AH_A>0$ and hence $\Delta_A^\ell>0$.

For the second family, the large-$s$ branch is already positive.  With
$s=u+3$ and $u\geq0$, direct expansion gives
\begin{align}
L_C={}&2q^2u^2+14q^2u+8q^2+6qtu^2+58qtu+120qt\notag\\
&+qu^4+22qu^3+177qu^2+616qu+780q\notag\\
&+4t^2u^2+44t^2u+112t^2+7tu^3+111tu^2\notag\\
&+584tu+1020t+2u^4+45u^3+379u^2+1416u+1980,
\label{eq:Clarge}
\end{align}
which is strictly positive.  The two remaining values of $s$ have the exact
certificates
\begin{align}
L_C+\frac{3H_C}{4t+11}
&=\frac{2(71qt^2+403qt+552q+62t^3+514t^2+1356t+1128)}{4t+11},
&&s=1,\label{eq:Csmall1}\\
L_C+\frac{H_C}{5t+22}
&=\frac{349qt^2+3138qt+7025q+338t^3+3976t^2+14950t+17600}{5t+22},
&&s=2.\label{eq:Csmall2}
\end{align}
All coefficients and denominators in Eqs.~\eqref{eq:Csmall1}--\eqref{eq:Csmall2}
are positive.  For $s\geq3$, Eq.~\eqref{eq:Clarge} directly gives
$\Delta_C^\ell>0$.  For $s=1,2$, use the high-core gap if $H_C>0$; otherwise
Eqs.~\eqref{eq:Csmall1}--\eqref{eq:Csmall2} imply $L_C>0$, so the low-core
gap is positive.  This proves the strict first inequality in Eq.~\eqref{eq:main}.
Finally, $s+1$ and $s+t+1$ are legal integer pineapple parameters at the same
$n$, proving the second inequality.  \hfill$\square$

The proof is computer-assisted only at the level of exact rational expansion
and identity checking.  It contains no floating-point sign decision, finite
parameter cutoff, interpolation, or asymptotic extrapolation.  The accompanying
derivation script reconstructs Eqs.~\eqref{eq:Agaps}--\eqref{eq:Csmall2} from
Eq.~\eqref{eq:sourceformula} and asserts the identities symbolically.  A second
implementation evaluates explicit code strings and selected literal full-graph
fundamental matrices.  Those finite checks protect against transcription and
indexing errors; they are not substitutes for the all-parameter identities.

\section{Why filling the core is not monotone}
\label{sec:nonmonotone}

When $t=1$, filling the interior zero turns either split-core graph into the
pineapple with parameter $r=s+2$.  The sign of the resulting change is not
fixed.  For family $A$, it is the sign of
\begin{align}
N_A(s,q)={}&(4s^2+16s+10)q^2+(s^3+9s^2+18s+4)q\notag\\
&-(s^5+13s^4+66s^3+163s^2+193s+84),
\label{eq:NA}
\end{align}
because the exact filled-minus-defect difference is
\begin{equation}
\frac{4N_A(s,q)}{(s+1)(s+3)(s+4)
(2q+s^2+7s+8)(2q+s^2+7s+12)}.
\label{eq:DAfill}
\end{equation}
For family $C$, the corresponding sign polynomial is
\begin{align}
N_C(s,q)={}&(6s+24)q^2+(-s^3-5s^2+2s+36)q\notag\\
&-2(s+3)(s+4)(s^2+6s+6),
\label{eq:NC}
\end{align}
and the difference is $(s+1)N_C$ divided by
\begin{equation}
(s+3)(s+4)(2q+s^2+5s+10)(2q+s^2+7s+12).
\end{equation}
Each sign polynomial is an upward-opening quadratic in $q$ with negative
constant term and hence has exactly one positive real root.  The change is
negative below that root and positive above it, with an integer equality
possible at the root.

For example,
\begin{equation}
 \K(011101)-\K(010101)
 =\frac{489}{110}-\frac{821}{180}=-\frac{229}{1980},
\end{equation}
whereas $s=1,q=4$ in the same family gives a positive change $11/840$.
In the second family, $s=12,q=48$ gives exact equality, with both values
$1055/16$.  The alternative pineapple with $r=13$ is nevertheless larger by
$337/2640$, as Theorem~1 requires.

This reversal is the transport content behind the two-pineapple proof.  Filling
the core adds edges, but it also changes local escape probabilities and the
degree-stationary target distribution.  An argument that compares only the
edge sets therefore misses the branch change in Eqs.~\eqref{eq:Agaps} and
\eqref{eq:Cgaps}.

\section{Relation to prior work and limitations}
\label{sec:related}

The closest result is the construction-code formula and pineapple conjecture of
Breen et al.~\cite{BreenEtAl2025}.  Their paper supplies
Eqs.~\eqref{eq:sourceformula} and \eqref{eq:pineapple}, but the checked preprint
does not state the two-family envelope proved here.  Recent work on Braess
cliques develops quotient and twin-vertex methods for changes in Kemeny's
constant~\cite{BreenDeBlieckVanderMeulen2026}; the checked results do not give
the present split-core comparison.  Effective-resistance matrices of weighted
threshold graphs have also been obtained~\cite{KeVanMieghem2026}, but that work
uses the combinatorial Laplacian and does not state this random-walk extremum.
Accordingly, novelty is claimed only within the primary texts actually checked,
not as an exhaustive priority certificate.

Theorem~1 has several sharp scope limits.  It covers exactly the two placements
in Eq.~\eqref{eq:families}, with every run nonempty.  It does not cover the other
six-run arrangements, multiple separated zero blocks, or an arbitrary threshold
code.  It neither classifies equality cases of the global problem nor proves
that every non-pineapple is strictly suboptimal.  The full pineapple conjecture
therefore remains unresolved.  The model is an exact finite, unweighted,
undirected, non-lazy walk; no continuum, thermodynamic, experimental, or
perturbative claim is implied.

\section{Conclusion}

Two infinite split-core families obey the conjectured pineapple bound for all
positive run lengths, and in fact each graph is strictly dominated by one of
two explicitly identified pineapples of the same order.  The need for a maximum
of two comparators is structural: direct core filling has all three possible
signs.  A natural next step is to seek a finite pineapple envelope for the
remaining six-run placements and then for two separated interior zero blocks.
Such an extension requires new identities; the present theorem alone does not
resolve the global extremal question.

\section*{Data and code availability}

The exact derivation script, independent checker, machine-readable polynomial
certificate, selected exact replay records, and reproduction instructions are
included with the source package.  No empirical dataset or stochastic sampling
is used.

\section*{Declaration of generative AI and AI-assisted technologies in the
manuscript preparation process}

OpenAI Codex was used to assist with problem development, symbolic code
construction, exact calculations, computational checking, literature
organization, and manuscript drafting.  Its outputs remain subject to review,
correction, and approval by the named human author before any submission.  The
named author must take responsibility for the content of any version that is
submitted; this local draft does not represent that approval.

\bibliographystyle{plainnat}
\bibliography{references}

\end{document}
